\chapter{Tenue de marché d'options en pratique}\label{ch:dv:options-market-making-in-practice}

La veille d'une date de détachement, deux teneurs de marché échangent entre eux 33\,000 contrats d'un call profondément dans la monnaie, dans les deux sens.
Aucun des deux ne prend de position. Tous deux exercent chaque call qu'ils ont acheté. Le lendemain matin, chacun constate que certains des calls qu'il a vendus n'ont jamais été assignés~:
leurs détenteurs, parmi le public, ne les ont pas exercés. Chaque call vendu non assigné vaut le dividende diminué d'une faible valeur temps. Dans le
modèle de ce chapitre, avec 30~\% des 10\,000 contrats du public non exercés, la paire dégage 108\,553~\$ après frais. Le \omterm{def:dv:options-market-making-in-practice:play}{dividend play} n'est qu'un petit coin du
market making d'options, mais il montre le caractère du métier~: de nombreux petits edges bien définis, chacun conçu avec soin, chacun avec ses propres coûts et risques. Ce
chapitre construit les autres~: une surface en temps réel et des \omterm{def:dv:options-market-making-in-practice:theo}{valeurs théoriques}, des largeurs de cotation, des politiques de couverture en delta, les dividendes et le risque de pin, et les
limites qui maintiennent le risque du portefeuille dans des buckets.

\section{Ajustement d'une surface en temps réel}

Un teneur de marché cote des milliers de séries. Il ne peut pas évaluer chacune à partir de zéro à mesure que le sous-jacent bouge, donc il conserve une surface~: un petit ensemble de
paramètres par échéance (le SVI du chapitre~8, ou un modèle de skew), recalibrés à intervalles réguliers sur les cotations et les transactions du marché, et déplacés avec le spot
entre-temps par une règle de marquage (chapitre~7).

\begin{definition}[Valeur théorique]\label{def:dv:options-market-making-in-practice:theo}
La \emph{valeur théorique}\index{valeur théorique} d'une option, pour un teneur de marché, est son prix sur la surface courante du teneur de marché,
avec sa courbe forward et sa courbe d'actualisation~: la référence autour de laquelle il fixe son cours acheteur et son cours vendeur, mise à jour à chaque mouvement du sous-jacent et à chaque recalibrage.
\end{definition}

Entre deux recalibrages, la règle de marquage fait le travail. Si le smile du marché bouge avec le spot (\omterm{def:dv:implied-volatility-and-its-surface:sticky}{sticky delta}) mais que la \omterm{def:dv:options-market-making-in-practice:theo}{valeur théorique} conserve la volatilité de chaque strike
(\omterm{def:dv:implied-volatility-and-its-surface:sticky}{sticky strike}), la \omterm{def:dv:options-market-making-in-practice:theo}{valeur théorique} dérive d'un terme égal au skew multiplié par la variation logarithmique du spot depuis le recalibrage, divisée par $\sqrt\tau$. Avec un skew de $-0.10$ et un
spot à 20~\% de volatilité, l'erreur moyenne en points de volatilité est~:

\begin{center}\small
\begin{tabular}{lrrrr}
\toprule
intervalle de recalibrage & 1 seconde & 10 secondes & 1 minute & 5 minutes\\
\midrule
options à un mois & 0.0023 & 0.0072 & 0.018 & 0.039\\
options à trois mois & 0.0013 & 0.0042 & 0.010 & 0.023\\
\bottomrule
\end{tabular}
\end{center}

Une minute de retard coûte un cinquantième de point de volatilité sur une option à un mois, largement en deçà de n'importe quelle largeur. Ce qui fait dérailler les surfaces n'est pas
l'intervalle mais les sauts~: une publication de résultats, une grosse transaction ou un mouvement d'indice déplace tout le smile d'un coup. Le moteur de cotation recalibre donc sur
les événements autant qu'à intervalles réguliers, et retire ses cotations pendant le recalibrage.

\section{Valeur théorique et largeurs de cotation}

Autour de la \omterm{def:dv:options-market-making-in-practice:theo}{valeur théorique}, le teneur de marché cote une largeur. Trop étroite, et les traders informés la saisissent dès que sa \omterm{def:dv:options-market-making-in-practice:theo}{valeur théorique} est fausse~; trop large, et le
flux non informé qui paie pour tout va ailleurs. Le livre~2 (chapitre~30) traitait la largeur dans un jeu~; la construction en fait une formule par
série.

\begin{example}[Le modèle de largeur]\label{ex:dv:options-market-making-in-practice:width}
On mesure tout en points de volatilité. L'erreur de la \omterm{def:dv:options-market-making-in-practice:theo}{valeur théorique} suit une loi normale d'écart-type 0.5 point. Les ordres non informés arrivent à raison de 20 par heure
lorsque la demi-largeur est nulle, décroissant comme $e^{-w/0.3}$ à mesure qu'elle s'élargit~; chacun paie la demi-largeur $w$. Les ordres informés arrivent à $\lambda$ par heure et
n'interviennent que lorsque l'erreur dépasse $w$, en prenant l'excédent. Le profit espéré par heure, par unité de \omterm{def:dv:greeks-and-the-hedging-pnl:greeks}{vega}, est
\[
20\,e^{-w/0.3}\,w-2\lambda\bigl(0.5\,\varphi(w/0.5)-w\,(1-\Phi(w/0.5))\bigr).
\]
Sans flux informé, la meilleure demi-largeur est de 0.30 point et rapporte 2.21 par heure. Avec 2, 5, 10 et 20 ordres informés par heure, elle s'élargit à 0.35,
0.42, 0.55 et 0.78, et le profit tombe à 1.89, 1.51, 1.07 et 0.65
(voir la \cref{fig:dv:options-market-making-in-practice:width}).
\end{example}

\begin{omfigure}
\begin{tikzpicture}
\begin{axis}[width=10.5cm,height=5cm,xlabel={demi-largeur (points de volatilité)},ylabel={profit par heure (par unité de vega)},
  tick label style={font=\small},label style={font=\small},
  xmin=0,xmax=1.5,ymin=-1,ymax=2.5,grid=major,grid style={gray!25},
  legend columns=5,legend cell align=left,legend style={font=\footnotesize,at={(0.5,-0.3)},anchor=north,draw=none,fill=none,/tikz/every even column/.append style={column sep=6pt}},
  table/col sep=comma]
\addplot[omDef,very thick] table[x=w,y=p0]{figdata/derivatives/26-options-market-making-in-practice/widths.csv};
\addplot[omThm,thick] table[x=w,y=p2]{figdata/derivatives/26-options-market-making-in-practice/widths.csv};
\addplot[omProp,thick] table[x=w,y=p5]{figdata/derivatives/26-options-market-making-in-practice/widths.csv};
\addplot[violet,thick] table[x=w,y=p10]{figdata/derivatives/26-options-market-making-in-practice/widths.csv};
\addplot[black,thick,dashed] table[x=w,y=p20]{figdata/derivatives/26-options-market-making-in-practice/widths.csv};
\addplot[black,only marks,mark=*,mark size=1.6pt] table[x=w,y=profit]{figdata/derivatives/26-options-market-making-in-practice/widths_opt.csv};
\legend{sans flux informé,2 par heure,5,10,20}
\end{axis}
\end{tikzpicture}

\omcaption{Profit espéré par heure de la cotation d'une demi-largeur autour de la \omterm{def:dv:options-market-making-in-practice:theo}{valeur théorique}, par unité de \omterm{def:dv:greeks-and-the-hedging-pnl:greeks}{vega}, pour plusieurs taux d'ordres informés~; les points marquent la
meilleure demi-largeur. Le flux informé élargit la largeur et fait baisser le profit. Données~: le tutoriel.}\label{fig:dv:options-market-making-in-practice:width}
\end{omfigure}

Exprimée en points de volatilité, la meilleure largeur ne dépend pas de l'option. En prix, elle vaut cette largeur multipliée par le \omterm{def:dv:greeks-and-the-hedging-pnl:greeks}{vega} de l'option. Une option à un mois
à la monnaie sur une action à 100 a un \omterm{def:dv:greeks-and-the-hedging-pnl:greeks}{vega} de 0.115 par point, contre 0.199 pour une option à trois mois, de sorte qu'à 0.42 point les demi-largeurs valent 4.8 et 8.4 centimes.
Les données d'entrée sont mesurées, non devinées~: l'erreur de la \omterm{def:dv:options-market-making-in-practice:theo}{valeur théorique}, mesurée à partir des mark-outs (livre~2, chapitre~15) des transactions passées comparées aux \omterm{def:dv:options-market-making-in-practice:theo}{valeurs théoriques} ultérieures, et la part
informée, tirée des mark-outs par contrepartie.

\section{Politiques de couverture en delta}

Le portefeuille du teneur de marché accumule du gamma à partir des transactions des clients, et son delta bouge à chaque tick. Couvrir chaque mouvement coûte la fourchette du
sous-jacent~; ne couvrir que rarement laisse subsister du risque. La couverture périodique, à intervalles réguliers, ignore où se situe le delta. Une \omterm{def:dv:options-market-making-in-practice:band}{bande de couverture} n'intervient que lorsqu'il le faut.

\begin{definition}[Bande de couverture]\label{def:dv:options-market-making-in-practice:band}
Une \emph{bande de couverture}\index{bande de couverture} est un intervalle autour du delta cible à l'intérieur duquel la couverture est laissée inchangée~; lorsque le delta en sort, la
couverture se recale sur le bord le plus proche de la bande, non sur son centre, si bien que les transactions ne sont jamais plus grandes que ce qu'exige le risque.
\end{definition}

La largeur de la bande découle d'un compromis. Le delta diffuse à un rythme fixé par le gamma, $\Gamma S\sigma$ par unité de $\sqrt t$, de sorte qu'il quitte une bande de
demi-largeur $H$ au bout d'un temps de l'ordre de $H^2/(\Gamma S\sigma)^2$, et coûte environ $\lambda SH$ en fourchette à chaque fois, avec $\lambda$ le coût proportionnel.
La variance de la partie non couverte croît comme $H^2S^2\sigma^2$. En minimisant le taux de coût plus l'aversion au risque $\gamma$ multipliée par le taux de variance, on obtient
$H^3\propto\lambda S\Gamma^2/\gamma$~: une demi-largeur de bande qui varie comme la racine cubique du coût multiplié par le carré du gamma. Whalley et Wilmott ont dérivé la bande optimale
pour de faibles coûts, à partir d'un modèle de maximisation de l'utilité, sous forme d'une expression du gamma de l'option. La construction utilise $H=c\,(\lambda S\Gamma^2)^{1/3}$ et traite
$c$ comme bouton de réglage du risque du desk.

\begin{example}[Bandes face à la couverture périodique]\label{ex:dv:options-market-making-in-practice:bands}
Un straddle vendu à un mois, à la monnaie, sur une action à 100 (volatilité de 20~\%), est couvert sur toute sa durée, avec un contrôle 13 fois par jour, à un coût de 5 points de base
du notionnel échangé. Couvrir à chaque contrôle coûte 0.525 en moyenne et laisse un écart-type du P\&L de 0.321. Une fois par jour coûte 0.142, avec un
écart-type de 0.855. Une bande avec $c=1$ coûte 0.140, comme la couverture quotidienne, avec un écart-type de 0.505, soit 41~\% de moins. Avec $c=0.25$, elle coûte 0.313
avec un écart-type de 0.313, et la couverture à chaque troisième contrôle coûte à peu près autant (0.302), avec un écart-type de 0.434
(voir la \cref{fig:dv:options-market-making-in-practice:bands}).
\end{example}

\begin{omfigure}
\begin{tikzpicture}
\begin{axis}[width=10.5cm,height=5.2cm,xlabel={coût moyen de couverture},ylabel={écart-type du P\&L},
  tick label style={font=\small},label style={font=\small},
  xmin=0,xmax=0.6,ymin=0,ymax=1.4,grid=major,grid style={gray!25},
  legend columns=2,legend cell align=left,legend style={font=\footnotesize,at={(0.5,-0.3)},anchor=north,draw=none,fill=none,/tikz/every even column/.append style={column sep=8pt}},
  table/col sep=comma]
\addplot[omProp,very thick,mark=*,mark size=1.6pt] table[x=cost,y=sd]{figdata/derivatives/26-options-market-making-in-practice/hedge_time.csv};
\addplot[omThm,very thick,mark=square*,mark size=1.6pt] table[x=cost,y=sd]{figdata/derivatives/26-options-market-making-in-practice/hedge_band.csv};
\legend{couverture périodique (tous les 1 à 26 contrôles),bande de couverture ($c$ de 0.25 à 4)}
\end{axis}
\end{tikzpicture}

\omcaption{Coût de couverture face au risque résiduel pour un straddle court d'un mois~: couverture périodique à intervalles de 1 à 26 contrôles (13 contrôles par jour),
et bandes de largeur croissante. La frontière des bandes se situe sous celle de la couverture périodique. Données~: le tutoriel.}\label{fig:dv:options-market-making-in-practice:bands}
\end{omfigure}

\section{Risque de pin, dividendes et exercice anticipé}

Certains risques du teneur de marché n'apparaissent dans aucune grecque. Le livre~1 (chapitre~23) a décrit le risque de pin. Un desk vendeur de 1\,000 calls de strike 100, avec une action qui clôture
à 100.02 le jour de l'échéance, ne sait pas combien seront exercés, entre aucun et la totalité, de sorte qu'il ignore s'il ouvrira le lundi
une position vendeuse de 100\,000 actions ou à plat. Il clôture ces positions avant la cloche ou détient la position en actions qu'implique l'assignation espérée, et accepte le reste.

Les dividendes créent l'opportunité inverse. La veille d'une date de détachement, les détenteurs d'options d'achat américaines profondément dans la monnaie devraient exercer lorsque le dividende
dépasse la valeur temps qu'ils abandonnent, soit la valeur de l'option de vente de même strike et même échéance (chapitre~6). Beaucoup ne le font pas. Pool, Stoll et Whaley ont établi
qu'entre 1996 et 2006, plus de la moitié des positions longues qui auraient dû être exercées ne l'ont pas été, pour un coût supérieur à 491~millions de dollars pour les détenteurs, et que les teneurs de marché
en ont capté l'essentiel.

\begin{definition}[Dividend play]\label{def:dv:options-market-making-in-practice:play}
Un \emph{dividend play}\index{dividend play} est une transaction réalisée la veille d'une date de détachement, au cours de laquelle les teneurs de marché achètent et vendent entre eux de grandes quantités égales et opposées d'une option d'achat
profondément dans la monnaie et exercent toutes leurs options d'achat longues~; comme l'assignation touche proportionnellement toutes les positions courtes, leurs calls vendus absorbent
l'essentiel de la non-assignation laissée par les détenteurs qui n'exercent pas, et chaque call vendu non assigné rapporte le dividende diminué de la valeur temps.
\end{definition}

Le calcul est simple. Avec $N$ contrats publics longs dont une part $f$ n'est pas exercée, et $q$ contrats échangés deux fois entre deux teneurs de marché,
$2q+(1-f)N$ calls sont exercés face à $N+2q$ positions courtes. Les calls vendus non assignés des teneurs de marché sont au nombre de $2q\,fN/(N+2q)$, qui tend vers $fN$ lorsque
$q$ croît. Les frais croissent linéairement en $q$, il existe donc une taille optimale, $q^*=\bigl(\sqrt{2fNg\,N/c}-N\bigr)/2$, où $g$ désigne le gain par call non assigné et $c$ les frais
par $q$.

\begin{example}[Une série de calls]\label{ex:dv:options-market-making-in-practice:play}
Un dividende de 50~cents (50~\$ par contrat), une option de vente valant 2~cents (2~\$), soit un gain de 48~\$ par contrat non assigné~; 10\,000 contrats publics~; frais (hypothèse)
de 0.10~\$ par contrat et par transaction et de 0.05~\$ par exercice, soit 0.50~\$ par $q$ pour quatre transactions et deux exercices. Si 30~\% des détenteurs n'exercent pas, la meilleure transaction
porte sur 32\,947 contrats. Elle laisse 2\,605 calls vendus non assignés, soit 87~\% des 3\,000 qui n'ont pas été exercés, pour un gain brut de 125\,026~\$, des frais de 16\,474~\$ et un net
de 108\,553~\$. Avec 10~\% de non-exercice, la meilleure transaction porte sur 16\,909 contrats et rapporte un net de 28\,591~\$~; avec 50~\%, 43\,990 contrats et un net de 193\,510~\$
(\cref{fig:dv:options-market-making-in-practice:play}).
\end{example}

\begin{omfigure}
\begin{tikzpicture}
\begin{axis}[width=10.5cm,height=5cm,xlabel={contrats échangés dans chaque sens (milliers)},ylabel={profit net espéré (milliers de \$)},
  tick label style={font=\small},label style={font=\small},
  xmin=0,xmax=100,ymin=-20,ymax=210,grid=major,grid style={gray!25},
  legend columns=3,legend cell align=left,legend style={font=\footnotesize,at={(0.5,-0.3)},anchor=north,draw=none,fill=none,/tikz/every even column/.append style={column sep=8pt}},
  table/col sep=comma]
\addplot[omDef,thick] table[x=q,y=f10]{figdata/derivatives/26-options-market-making-in-practice/dividend.csv};
\addplot[omProp,very thick] table[x=q,y=f30]{figdata/derivatives/26-options-market-making-in-practice/dividend.csv};
\addplot[omThm,thick] table[x=q,y=f50]{figdata/derivatives/26-options-market-making-in-practice/dividend.csv};
\addplot[black,only marks,mark=*,mark size=1.6pt] table[x=q,y=net]{figdata/derivatives/26-options-market-making-in-practice/dividend_opt.csv};
\legend{10~\% de non-exercice,30~\%,50~\%}
\end{axis}
\end{tikzpicture}

\omcaption{Profit net espéré d'un \omterm{def:dv:options-market-making-in-practice:play}{dividend play} sur une série de calls (10\,000 contrats publics, gain de 48~\$ par call non assigné, 0.50~\$ de frais par contrat échangé
dans chaque sens) en fonction de la taille de la transaction, pour trois parts de détenteurs n'exerçant pas~; les points marquent la taille optimale. Données~: le
tutoriel.}\label{fig:dv:options-market-making-in-practice:play}
\end{omfigure}

Le profit espéré n'est pas le profit réalisé~: l'assignation est aléatoire, et un desk assigné au-delà de sa part conserve moins. La transaction existe parce que
les frais et les règles d'assignation la rendent peu coûteuse, et elle disparaît là où ils ne le font pas.

\begin{dated}{2026-09}{Plafonds de frais de stratégie}\label{dat:dv:options-market-making-in-practice:fees}
Le barème de frais de Cboe Options du 15 septembre 2026 plafonne à 0.00~\$ les frais de transaction des teneurs de marché et des autres non-clients pour les stratégies de merger, de short-stock-interest, de reversal,
de conversion et de jelly-roll exécutées en open outcry le même jour dans la même classe. Les stratégies de dividende ne figurent pas sur cette liste, de sorte qu'un \omterm{def:dv:options-market-making-in-practice:play}{dividend play}
y paie les frais ordinaires.
\end{dated}

\section{Limites par bucket}

Le \omterm{def:dv:greeks-and-the-hedging-pnl:greeks}{vega} d'un teneur de marché n'est pas un nombre unique. Un \omterm{def:dv:greeks-and-the-hedging-pnl:greeks}{vega} long d'un mois et un \omterm{def:dv:greeks-and-the-hedging-pnl:greeks}{vega} court d'un an de même taille sont neutres en total et entièrement exposés à une variation de la
structure par terme. Les limites se fixent par bucket.

\begin{definition}[Bucket de vega]\label{def:dv:options-market-making-in-practice:bucket}
Un \emph{bucket de vega}\index{bucket de vega} est une plage d'échéances (et, dans les schémas plus fins, de moneyness) sur laquelle le \omterm{def:dv:greeks-and-the-hedging-pnl:greeks}{vega} d'un livre est sommé et limité~; le profil de \omterm{def:dv:greeks-and-the-hedging-pnl:greeks}{vega} d'un livre
est son \omterm{def:dv:greeks-and-the-hedging-pnl:greeks}{vega} par bucket, et les limites s'appliquent séparément à chaque bucket.
\end{definition}

Lorsqu'un bucket approche de sa limite, le moteur de cotation ne s'arrête pas aussitôt. Il décale les deux cotations dans le sens qui attire des transactions de compensation~: un desk
vendeur de \omterm{def:dv:greeks-and-the-hedging-pnl:greeks}{vega} relève son cours acheteur et son cours vendeur afin que les clients lui vendent des options. Il retire le côté qui franchirait la limite seulement lorsque celle-ci est atteinte.

\begin{example}[Un mois de demande de puts]\label{ex:dv:options-market-making-in-practice:toy}
Un teneur de marché jouet cote 15 séries (un, trois et six mois~; strikes de 90 à 110) pendant 20 jours, à raison de 13 contrôles par jour. Des ordres de 10 contrats arrivent trois fois par
contrôle~; les clients non informés achètent des puts deux fois sur trois et des calls la moitié du temps~; un ordre sur cinq est informé. Le theo se trompe de 0.5 point, la demi-fourchette vaut
0.4 point, et le delta est couvert à chaque contrôle à 5~pb. Les limites sont de 3\,000~\$, 5\,000~\$ et 6\,000~\$ de \omterm{def:dv:greeks-and-the-hedging-pnl:greeks}{vega} par point. Avec décalage, le \omterm{def:dv:greeks-and-the-hedging-pnl:greeks}{vega} du livre à six mois
culmine à $-3\,369$~\$ et le P\&L vaut 11\,141~\$~: edge 36\,280~\$, sélection $-8\,720$~\$, couverture $-9\,815$~\$, inventaire $-6\,605$~\$. Sans décalage, le même flux amène le
\omterm{def:dv:greeks-and-the-hedging-pnl:greeks}{vega} à six mois à $-5\,033$~\$, l'inventaire perd 15\,240~\$, et le mois se termine à $-607$~\$ (\cref{fig:dv:options-market-making-in-practice:vega}). Un seul parcours ne prouve
rien sur les moyennes. Il montre où va le risque lorsque rien ne s'oppose au flux.
\end{example}

\begin{omfigure}
\begin{tikzpicture}
\begin{groupplot}[group style={group size=2 by 1,horizontal sep=1.3cm},
  width=6.6cm,height=5cm,xlabel={contrôle (13 par jour)},xmin=0,xmax=260,ymin=-6.5,ymax=2.5,
  tick label style={font=\small},label style={font=\small},grid=major,grid style={gray!25},table/col sep=comma]
\nextgroupplot[ylabel={vega (milliers de \$ par point)},title={\small bucket à trois mois},
  legend columns=2,legend cell align=left,legend style={font=\footnotesize,at={(1.08,-0.3)},anchor=north,draw=none,fill=none,/tikz/every even column/.append style={column sep=8pt}}]
\addplot[omThm,very thick] table[x=step,y=m3_on]{figdata/derivatives/26-options-market-making-in-practice/vega.csv};
\addplot[omProp,thick] table[x=step,y=m3_off]{figdata/derivatives/26-options-market-making-in-practice/vega.csv};
\addplot[gray,dashed,domain=0:260] {-5};
\legend{cotations décalées par les limites,sans limites}
\nextgroupplot[title={\small bucket à six mois},yticklabels={}]
\addplot[omThm,very thick] table[x=step,y=m6_on]{figdata/derivatives/26-options-market-making-in-practice/vega.csv};
\addplot[omProp,thick] table[x=step,y=m6_off]{figdata/derivatives/26-options-market-making-in-practice/vega.csv};
\addplot[gray,dashed,domain=0:260] {-6};
\end{groupplot}
\end{tikzpicture}

\omcaption{\omterm{def:dv:greeks-and-the-hedging-pnl:greeks}{Vega} par bucket d'un teneur de marché d'options jouet face à un mois de demande de puts, avec des cotations décalées à mesure que chaque bucket approche de sa limite (pointillés) et
sans décalage. Données~: le tutoriel.}\label{fig:dv:options-market-making-in-practice:vega}
\end{omfigure}

Le gamma obéit au même traitement par échéance, et les plus gros noms reçoivent leurs propres buckets. Les limites ne valent que ce que valent les \omterm{def:dv:greeks-and-the-hedging-pnl:greeks}{grecques} qui les sous-tendent, c'est pourquoi
les \omterm{def:dv:greeks-and-the-hedging-pnl:greeks}{grecques} du desk proviennent de la même surface que ses cotations.

\section{Tutoriel~: un teneur de marché d'options jouet}\label{tut:dv:options-market-making-in-practice:tut}

\begin{tutorial}
\textbf{Objectif.} Calculer les meilleures demi-largeurs du modèle de largeur, comparer les \omterm{def:dv:options-market-making-in-practice:band}{bandes de couverture} à une couverture à intervalles fixes, dimensionner un \omterm{def:dv:options-market-making-in-practice:play}{dividend play} et faire tourner un teneur de marché d'options jouet avec
des limites par bucket. \textbf{État final~:} les quatre figures, le tableau de péremption et les chiffres du problème du week-end.
\begin{tutsteps}
\item \textbf{Le hedger}, à bande ou à horloge, vectorisé sur les trajectoires~:
\omcode{code/firm/optmm/firm_optmm.py}{77}{108}{Couverture en delta à intervalles fixes ou dans une bande.}\label{lst:dv:options-market-making-in-practice:hedge}
\item \textbf{Le \omterm{def:dv:options-market-making-in-practice:play}{dividend play}} et sa taille optimale~:
\omcode{code/firm/optmm/firm_optmm.py}{118}{137}{L'espérance de gain du dividend play.}\label{lst:dv:options-market-making-in-practice:play}
\item \textbf{La cotation}, décalée selon l'usage que fait le bucket de sa limite~:
\omcode{code/firm/optmm/firm_optmm.py}{160}{174}{Cotations autour de la valeur théorique, décalées selon la limite.}\label{lst:dv:options-market-making-in-practice:quote}
\item \textbf{Lancer} \texttt{dv\_optmm.stale\_table()}, \texttt{width\_study()}, \texttt{hedge\_study()}, \texttt{dividend\_study()}, \texttt{toy\_mm()} et
\texttt{fig\_optmm.py}.
\end{tutsteps}
\textbf{Pour aller plus loin.} Faire dépendre la demi-largeur de la part de flux informés par série~; couvrir le teneur de marché jouet avec une bande~; tirer au sort
l'assignation du \omterm{def:dv:options-market-making-in-practice:play}{dividend play} et tracer la distribution de son gain.
\end{tutorial}

\section{Construction~: le moteur de cotation d'options}\label{bld:dv:options-market-making-in-practice:optmm}

\begin{build}
\bfield{Rôle} Le moteur de cotation d'options de la petite société~: \omterm{def:dv:options-market-making-in-practice:theo}{valeurs théoriques} tirées d'une surface en direct, largeurs issues d'un modèle, un hedger à bande, l'arithmétique du \omterm{def:dv:options-market-making-in-practice:play}{dividend play},
et des limites par bucket qui décalent les cotations.
\bfield{Interface} \texttt{theo(option, spot, t, surface)} $\to$ (valeur, \omterm{def:dv:greeks-and-the-hedging-pnl:greeks}{vega} par point)~; \texttt{width\_profit}, \texttt{optimal\_width(err, lam\_u, w\_scale,
lam\_i)}~; \texttt{band\_half\_width}, \texttt{hedge\_paths(paths, book, vol, dt, cost, band, every)}~; \texttt{should\_exercise}, \texttt{dividend\_play(public,
fail, q, gain, trade\_fee, exercise\_fee)}, \texttt{best\_play}~; \texttt{Bucket(name, lo, hi, limit)}, \texttt{vega\_by\_bucket}, \texttt{quote(option, spot, t,
surface, half\_width\_vol, bucket\_vega, limit)}. \bfield{Règles}
Les largeurs sont fixées en points de volatilité et converties par le \omterm{def:dv:greeks-and-the-hedging-pnl:greeks}{vega}~; les couvertures sont ramenées jusqu'au bord de la bande~; une limite décale les cotations avant de retirer un côté~; le mark-out de chaque
transaction est enregistré pour le modèle de largeur.
\bfield{Tests d'acceptation} \texttt{code/firm/optmm/tests/}~: sans flux informé, la meilleure largeur est l'échelle de la demande, et un flux informé l'élargit~; un \omterm{def:dv:options-market-making-in-practice:play}{dividend play} infini
capte chaque défaut d'exercice et la meilleure taille bat ses voisines~; les cotations sont symétriques à usage nul, décalées quand on est long de \omterm{def:dv:greeks-and-the-hedging-pnl:greeks}{vega}, et un côté est
retiré à la limite~; les bandes coûtent moins que la couverture à chaque contrôle et portent davantage de risque que des bandes étroites.
\bfield{Pour aller plus loin} Assignation aléatoire dans le \omterm{def:dv:options-market-making-in-practice:play}{dividend play}~; largeurs par série à partir des mark-out enregistrés~; limites de gamma par échéance.
\end{build}

\begin{omsources}
\item M.~Pool, H.~R.~Stoll and R.~E.~Whaley, ``Failure to exercise call options: an anomaly and a trading game'', \emph{Journal of Financial Markets} 11(1)
(2008) 1--35.
\item A.~E.~Whalley and P.~Wilmott, ``An asymptotic analysis of an optimal hedging model for option pricing with transaction costs'', \emph{Mathematical Finance}
7(3) (1997) 307--324.
\item Cboe Exchange, Inc., fees schedule, 15 septembre 2026, note~13.
\end{omsources}

\section{Exercices}

\begin{exercise}[$\star$]\label{exo:dv:options-market-making-in-practice:1}
Pourquoi la meilleure demi-largeur est-elle égale à l'échelle de la demande (0.3 point) lorsqu'il n'y a pas de flux informé~?
\end{exercise}

\begin{exercise}[$\star$]\label{exo:dv:options-market-making-in-practice:2}
Convertir une demi-largeur de 0.42 point de volatilité en centimes pour des options à la monnaie à un mois et à trois mois sur une action cotée à 100.
\end{exercise}

\begin{exercise}[$\star$]\label{exo:dv:options-market-making-in-practice:3}
Vérifier le décompte des calls non assignés du \omterm{def:dv:options-market-making-in-practice:play}{dividend play} pour $N=10\,000$, $f=0.3$, $q=32\,947$.
\end{exercise}

\begin{exercise}[$\star\star$]\label{exo:dv:options-market-making-in-practice:4}
Dériver la meilleure taille de transaction $q^*$ du \omterm{def:dv:options-market-making-in-practice:play}{dividend play}.
\end{exercise}

\begin{exercise}[$\star\star$]\label{exo:dv:options-market-making-in-practice:5}
Pourquoi une \omterm{def:dv:options-market-making-in-practice:band}{bande de couverture} s'exécute-t-elle jusqu'à son bord plutôt que jusqu'à son centre~?
\end{exercise}

\begin{exercise}[$\star\star$]\label{exo:dv:options-market-making-in-practice:6}
Pourquoi l'erreur liée à une surface périmée croît-elle comme la racine carrée de l'intervalle de recalibrage et décroît-elle avec l'échéance de l'option~?
\end{exercise}

\begin{exercise}[$\star\star\star$]\label{exo:dv:options-market-making-in-practice:7}
\emph{Programmation.} Tracer les assignations du \omterm{def:dv:options-market-making-in-practice:play}{dividend play} au hasard (chaque call court étant assigné avec la même probabilité) pour $f=0.3$ et la meilleure taille, et donner
l'écart type du gain net.
\end{exercise}

\begin{exercise}[$\star\star\star$]\label{exo:dv:options-market-making-in-practice:8}
\emph{Trouver le défaut.} «~Notre teneur de marché jouet a gagné 11\,141~\$ avec des limites et perdu 607~\$ sans~: les limites ajoutent environ 11\,700~\$ par mois.~»
\end{exercise}

\section{Problème~: La matinée du dividend play}

\begin{problem}[{Problème du week-end --- ce que valent ceux qui n'exercent pas}]\label{pb:dv:options-market-making-in-practice:1}
Une action verse un dividende de 50 cents demain. Sa série de calls profondément dans la monnaie compte 10\,000 contrats de positions ouvertes du public~; le put de même prix d'exercice vaut 2
cents. Deux teneurs de marché envisagent un \omterm{def:dv:options-market-making-in-practice:play}{dividend play}.

\medskip\noindent\textbf{Partie I --- Exercice.}
\begin{enumerate}
\item Le détenteur du call doit-il exercer ce soir, et que coûte par contrat le fait de ne pas exercer~?
\item Pourquoi l'exercice anticipé d'un call ne peut-il être optimal que juste avant une date de détachement (à taux nuls)~?
\item Que Pool, Stoll et Whaley ont-ils constaté sur la fréquence à laquelle les détenteurs omettent d'exercer~?
\item Qui détient l'autre côté des calls du public, et que leur arrive-t-il lorsque les détenteurs omettent d'exercer~?
\item Pourquoi l'assignation au prorata sur les positions courtes rend-elle le jeu rentable~?
\end{enumerate}

\medskip\noindent\textbf{Partie II --- Le jeu.}
\begin{enumerate}[resume]
\item Écrire le nombre de calls courts non assignés des teneurs de marché en fonction de $q$.
\item Quel est le gain maximal qu'ils peuvent capter, et pour quelle taille~?
\item Donner les frais par contrat échangé dans chaque sens selon l'hypothèse du chapitre.
\item Calculer la meilleure taille et le gain net lorsque 30~\% des détenteurs omettent d'exercer.
\item Donner le gain net lorsque 10~\% et lorsque 50~\% des détenteurs omettent d'exercer.
\end{enumerate}

\medskip\noindent\textbf{Partie III --- Risques.}
\begin{enumerate}[resume]
\item Que se passe-t-il si les teneurs de marché sont assignés plus que leur part au prorata~?
\item Que se passe-t-il si le public exerce davantage que prévu~?
\item Comment l'absence de plafond de frais pour les stratégies sur dividende modifie-t-elle la transaction~?
\item Quel autre risque le desk porte-t-il pendant la nuit sur ces positions~?
\item Comment estime-t-on à l'avance le taux de défaut d'une série~?
\end{enumerate}

\medskip\noindent\textbf{Partie IV --- Jugement.}
\begin{enumerate}[resume]
\item Qui perd dans ce jeu, et est-ce l'œuvre des teneurs de marché~?
\item Le desk doit-il réaliser la transaction avec un taux d'échec attendu de 10~\%~?
\item Comment une modification de règle sur l'exercice ou sur les frais mettrait-elle fin au jeu~?
\item Énoncer le \emph{résultat nommé}~: le profit espéré d'un \omterm{def:dv:options-market-making-in-practice:play}{dividend play} sur une série d'options d'achat lorsqu'une part donnée des détenteurs n'exerce pas, net des frais.
\item En une phrase~: d'où provient le profit du \omterm{def:dv:options-market-making-in-practice:play}{dividend play}~?
\end{enumerate}
\end{problem}

\section{Questions d'entretien}

\begin{interviewq}[$\star$ \iqroles{trader}]\label{iq:dv:options-market-making-in-practice:1}
Comment fixez-vous la largeur d'une cotation d'option~?
\end{interviewq}

\begin{interviewq}[$\star\star$ \iqroles{trader, chercheur}]\label{iq:dv:options-market-making-in-practice:2}
À quelle fréquence un teneur de marché d'options doit-il couvrir son delta~?
\end{interviewq}

\begin{interviewq}[$\star\star$ \iqroles{trader}]\label{iq:dv:options-market-making-in-practice:3}
Expliquez un \omterm{def:dv:options-market-making-in-practice:play}{dividend play}. Qui y perd~?
\end{interviewq}

\begin{interviewq}[$\star\star$ \iqroles{développeur}]\label{iq:dv:options-market-making-in-practice:4}
Concevez la partie d'un moteur de cotation qui met à jour les \omterm{def:dv:options-market-making-in-practice:theo}{valeurs théoriques} à chaque tick du sous-jacent. Qu'est-ce qui doit être rapide, et qu'est-ce qui peut être lent~?
\end{interviewq}

\begin{interviewq}[$\star\star$ \iqroles{risque, trader}]\label{iq:dv:options-market-making-in-practice:5}
Pourquoi limiter le \omterm{def:dv:greeks-and-the-hedging-pnl:greeks}{vega} par bucket plutôt qu'en total~? Que fait le desk lorsqu'un bucket approche de sa limite~?
\end{interviewq}

\begin{interviewq}[$\star\star\star$ \iqroles{trader}]\label{iq:dv:options-market-making-in-practice:6}
C'est le vendredi d'échéance et le titre se trouve sur votre plus grand strike vendu. Que faites-vous~?
\end{interviewq}
